% Figaro Paper A: Asymmetric Bonding and Buyer Dominance — Two Composing Mechanisms
% Mechanism-design paper. Self-contained.
% Preprint, May 2026

\documentclass[11pt,a4paper]{article}

% ── Packages ────────────────────────────────────────────────────────
\usepackage[utf8]{inputenc}
\usepackage[T1]{fontenc}
\usepackage{amsmath,amssymb,amsthm}
\usepackage{mathtools}
\usepackage{booktabs}
\usepackage{graphicx}
\usepackage{hyperref}
\usepackage{cleveref}
\usepackage{enumitem}
\usepackage{xcolor}
\usepackage[margin=1in]{geometry}
\usepackage{natbib}
\bibliographystyle{plainnat}

% ── Theorem environments ────────────────────────────────────────────
\newtheorem{theorem}{Theorem}[section]
\newtheorem{lemma}[theorem]{Lemma}
\newtheorem{corollary}[theorem]{Corollary}
\newtheorem{proposition}[theorem]{Proposition}
\theoremstyle{definition}
\newtheorem{definition}[theorem]{Definition}
\newtheorem{example}[theorem]{Example}
\theoremstyle{remark}
\newtheorem{remark}[theorem]{Remark}

% ── Macros ──────────────────────────────────────────────────────────
\newcommand{\figaro}{\textsc{Figaro}}
\newcommand{\pay}{P}           % payment
\newcommand{\cumval}{G}        % cumulative value
\newcommand{\sbond}{C_s}       % seller bond
\newcommand{\bbond}{C_b}       % buyer bond
\newcommand{\spay}{\pi_s}      % seller payout
\newcommand{\bpay}{\pi_b}      % buyer payout
\newcommand{\ub}{u_b}          % buyer utility
\newcommand{\us}{u_s}          % seller utility
\DeclareMathOperator{\Coop}{C}
\DeclareMathOperator{\Def}{D}

% ────────────────────────────────────────────────────────────────────

\title{%
  Asymmetric Bonding and Buyer Dominance:\\
  Two Composing Mechanisms for Self-Enforcing $N$-Party Coordination%
}

\author{
  Alessandro Daliana
}

\date{April 2026}

\begin{document}

\maketitle

% ════════════════════════════════════════════════════════════════════
\begin{abstract}
We present \figaro, a settlement mechanism that enables self-enforcing
agreements between mutually distrusting parties without arbitrators,
timeouts, or governance. The mechanism composes two distinct primitives.
\emph{Asymmetric bonding}: each party locks collateral equal to twice
the transaction value, making cooperation the weakly dominant strategy
in a one-shot game and the unique profile surviving iterated
elimination of weakly dominated strategies. \emph{Buyer dominance
with atomic resolution}: only the buyer can trigger resolution, and
resolution settles all orders in a process simultaneously or not at
all. The two mechanisms are inseparable: bonding alone gives
independently secured bilateral edges that cannot coordinate at the
multi-party level; buyer dominance alone is without force absent the
bonding equilibrium.

Asymmetric bonding scales naturally from two-party exchange to
$N$-party process chains: because each seller's bond is tied to the
cumulative value of all upstream work rather than to the local
payment, the Nash equilibrium is preserved at every chain position
and the seller's risk-to-reward ratio grows with chain depth. Buyer
dominance with atomic resolution then operates on the already-scaled
mesh: the all-or-nothing resolution rule induces a weakest-link
subgame among sellers. We show that the resulting peer-enforcement
properties obtain under strictly weaker assumptions than the
joint-liability mechanisms classically analyzed in the microfinance
literature.
\end{abstract}

\medskip
\noindent\textbf{Keywords:}
mechanism design, Nash equilibrium, collateral bonding,
multi-party coordination, process chains, peer enforcement

% ════════════════════════════════════════════════════════════════════
\section{Introduction}
\label{sec:intro}

\subsection{The Problem}

Economic coordination between strangers requires trust infrastructure.
When two parties who have never met wish to exchange value, they face
the classic commitment problem: either party may defect after the
other has performed. The historical solution is \emph{institutional
intermediation}---firms, courts, escrow services, platform
operators---each of which introduces costs, delays, jurisdictional
dependencies, and power asymmetries.

The problem intensifies in multi-party settings. A food delivery
involves a cook, a kitchen operator, an ingredient sourcer, and a
courier; a procurement process involves engineers, manufacturers,
inspectors, and logistics providers. The standard architectural
response is to aggregate participants into a single legal entity that
internalizes the coordination problem and resolves it through
hierarchical management.

Blockchain-based smart contracts have produced a rich literature on
trustless exchange, but most protocols address financial
operations---swaps, lending, derivatives---rather than general
coordination. Those that target coordination
(e.g.,~\citep{kleros2019,aragon2017}) typically reintroduce the
intermediary as an on-chain dispute resolution mechanism: a jury, a
governance body, or a timeout that unilaterally releases funds.

\subsection{Our Contribution}

We present \figaro{}, which achieves self-enforcing coordination
through two composing mechanisms: \emph{asymmetric bonding}, which
produces the bilateral Nash equilibrium and scales it naturally to
$N$-party process chains by tying each seller's bond to cumulative
upstream value; and \emph{buyer dominance with atomic resolution},
which enforces inter-seller coordination on the already-scaled mesh.
The two mechanisms are inseparable: bonding alone yields independently
secured edges that cannot coordinate at the multi-party level; buyer
dominance alone is worthless without the bonding equilibrium. Together
they make the mesh resolvable from a single signature with cooperation
pressure propagating through it, and they create the conditions for
participants to resolve performance among themselves as social agents:
the bonded commitment puts buyer and sellers in a position where
working out difficulties between them is the dominant strategy, and
external dispute resolution is a marginal residual rather than the
primary path (\Cref{sec:dispute}). The contributions are:

\begin{enumerate}[label=(\roman*)]
  \item \textbf{A minimal settlement primitive} with no timeout, no
    governance, and no escape hatches from the bonded state. We show
    that this minimality is not a simplification but a \emph{security
    requirement}: every additional exit path weakens the Nash
    equilibrium (\Cref{thm:escape-hatch-weakness}).

  \item \textbf{Asymmetric bonding scales to $N$-party chains.} We
    prove (\Cref{thm:n-party-nash}) that cooperation remains weakly
    dominant at every position in an arbitrary-length process chain
    (with strict dominance whenever at least one other player's
    cooperation is not already foreclosed), and is the unique profile
    surviving iterated elimination of weakly dominated strategies;
    the seller's risk-to-reward ratio increases with chain depth.

  \item \textbf{Atomic resolution} as a coordination forcing function.
    We show that all-or-nothing process settlement induces a
    weakest-link subgame among sellers, creating endogenous peer
    pressure without explicit communication or governance.

  \item \textbf{A reduction over joint-liability microfinance.} We
    prove (\Cref{prop:microfinance-reduction}) that the
    peer-enforcement outcome common to \citet{ghatak1999} and
    \citet{stiglitz1990} is obtained by the bonded commitment
    mechanism under strictly weaker assumptions: no repeated
    interaction, no local information, no exogenous punishment
    technology, and no joint-liability contracting.

\end{enumerate}

\subsection{Paper Organization}

\Cref{sec:related} surveys related work. \Cref{sec:mechanism} presents
the mechanism formally. \Cref{sec:game-theory} proves the
game-theoretic properties. \Cref{sec:process-trees} extends the model
to $N$-party process chains, including a comparison with
joint-liability microfinance (\Cref{sec:microfinance}).
\Cref{sec:dispute} treats dispute resolution as a marginal external
concern. \Cref{sec:conclusion} concludes.

\subsection{Architectural Posture}
\label{sec:architectural-posture}

\paragraph{The mechanism is ideologically agnostic; the graph is the
politics.} A consequence of locating the whole of the present
paper's argument at the mechanism level is that the mechanism admits no
commitments about \emph{what} is being coordinated or \emph{under
what normative regime}. Currency, jurisdiction, identity type,
arbitration forum, role structure, price-discovery mechanism,
contribution metric---none of these are in the mechanism; all of
them are expressed in the \emph{graph} that parties compose on top
of it. A Dutch-auction-and-private-operator graph expresses a
market-liberal coordination pattern; a
cooperative-bonding-and-contribution-weighted-resolution graph
expresses a cooperative one; an interest-free risk-sharing graph
expresses an Islamic-finance one. The mechanism enables each and
prefers none. This is not evasion. It is the architectural
consequence of placing the enforcement function at TCP/IP altitude:
a layer below which nothing remains to factor out, and a layer
above which the normative choices become legible as choices rather
than as properties of the infrastructure.

\paragraph{Composability is external-first.} A second consequence
of the mechanism's narrow construction is that its useful extensions
are predominantly \emph{external} rather than internal.
The mechanism does not include a dispute forum (Kleros, the Singapore
International Arbitration Centre, or any court); it does not
include a carbon-offset market (Klima DAO, Toucan, Moss); it does
not include a prediction market, an insurance pool, a lending
facility, a tax-reporting service, an identity provider, a storage
layer, or a messaging fabric. Each of these exists as a working
external protocol or institution, and an assembly composed on top
of the mechanism can name any of them in its graph: a forum-clause
names the dispute forum the parties agree to; an attestation
clause names the offset operator the buyer wants to compose;
a sub-order carries the buyer's offset purchase to the named
operator within the same atomic resolution as the original commerce.
The mechanism itself need not (and, by
\Cref{thm:escape-hatch-weakness}, must not) absorb the bookkeeping
that such composition implies. The composability posture is the
architectural counterpart of the ideological-agnosticism claim
above: the mechanism makes external surfaces composable; the graph
names which externals it composes with.

% ════════════════════════════════════════════════════════════════════
\section{Related Work}
\label{sec:related}

\paragraph{Mutual-destruction escrows.}
The direct ancestor of \figaro{} is the Safe Remote Purchase
contract~\citep{solidity2016}, included in the Solidity documentation
as a reference pattern: buyer and seller each lock $2\times$ the item
value, and only mutual agreement releases the funds. This two-party,
single-order pattern is the seed from which \figaro{} grew.
BitHalo~\citep{bithalo2014} independently explored mutual-destruction
deposits for peer-to-peer Bitcoin trade.

The central problem was generalizing SRP beyond two parties. SRP uses
symmetric $2\times$ deposits with mutual agreement to release---both
parties must cooperate. This works for a single exchange but cannot
compose: in a twelve-seller delivery workflow, there is no natural
``mutual agreement'' among all participants. \emph{Asymmetric bonding}
(\Cref{def:asymmetric-bonding}) solves this. By scaling each seller's
bond with cumulative upstream value rather than fixing it per trade,
the mechanism creates a mesh of independently secured edges. Each edge
carries its own Nash equilibrium (\Cref{thm:two-party-nash}), and the
mesh decomposes as finely as is operationally tolerable.

\emph{Buyer dominance} then takes advantage of this mesh structure to
make multi-party coordination resolvable from a single signature.
Because a single party resolves the entire process atomically,
coordination pressure propagates through the mesh without any
management layer (\Cref{thm:coordination-pressure}).

\paragraph{State channels and payment networks.}
The Lightning Network~\citep{poon2016} and Raiden~\citep{raiden2017}
use time-locked deposits for off-chain payment. These protocols
optimize for high-frequency bilateral payments and rely on timeouts
for dispute resolution---precisely the escape hatch that \figaro{}
eliminates (\Cref{thm:escape-hatch-weakness}). More critically, state
channels do not address multi-party \emph{coordination} (delivery,
fulfillment, service composition); they address multi-hop
\emph{payment routing}.

\paragraph{On-chain dispute resolution.}
Kleros~\citep{kleros2019} provides decentralized arbitration via
Schelling-point juror selection. Aragon Court~\citep{aragon2017}
offers similar functionality within DAO governance. Both reintroduce
a \emph{discretionary human decision} into the resolution path,
which breaks the self-enforcing property
(\Cref{thm:escape-hatch-weakness}, Case~2).

\paragraph{Optimistic mechanisms.}
Optimistic rollups~\citep{optimism2021} and optimistic
oracles~\citep{uma2020} use a challenge-response pattern: an assertion
stands unless challenged within a timeout window. This requires
liveness assumptions (the challenger must be online) and timeout
calibration (too short enables fraud; too long delays finality).
\figaro{} has no timeout from the bonded state---capital remains
locked indefinitely until the buyer resolves, making liveness
irrelevant to security.

\paragraph{Mechanism design.}
The revelation principle~\citep{myerson1979} guarantees that any
achievable outcome can be implemented via a direct mechanism.
\figaro{}'s insight is that the ``direct mechanism'' for multi-party
coordination need not involve message passing, reporting, or
arbitration at all: locked capital \emph{is} the mechanism. The
closest theoretical antecedent is the concept of \emph{deposit
games}~\citep{asgaonkar2019}, which analyze two-party security
deposits. We extend deposit-game analysis to $N$-party process chains
under asymmetric bonding.

\paragraph{Joint-liability microfinance.}
Group lending models---most prominently the Grameen
Bank~\citep{yunus1999} and its formalizations by
\citet{ghatak1999} and \citet{stiglitz1990}---show that joint
liability combined with a social substrate can produce peer
selection and peer monitoring outcomes. We return to these models in
\Cref{sec:microfinance}: \figaro{} reproduces the same
peer-enforcement outcome under strictly weaker assumptions, with no
appeal to repeated interaction or local information.

% ════════════════════════════════════════════════════════════════════
\section{The Two Mechanisms: Formal Definition}
\label{sec:mechanism}

\subsection{Notation and Setup}

Consider two parties, a \emph{buyer} $B$ and a \emph{seller} $S$, who
wish to exchange a payment $\pay > 0$ (denominated in a unit of
account) for goods or services. Both parties hold sufficient balances.

\begin{definition}[Asymmetric Bonding]
\label{def:asymmetric-bonding}
An \emph{asymmetric bond} for a transaction with payment $\pay$ and
cumulative value $\cumval \geq \pay$ is a pair of deposits:
\begin{align}
  \bbond &= 2\pay \quad (\text{buyer bond}), \label{eq:buyer-bond} \\
  \sbond &= 2\cumval \quad (\text{seller bond}). \label{eq:seller-bond}
\end{align}
For a root transaction (no upstream value), $\cumval = \pay$, so both
bonds equal $2\pay$. For sub-orders in a process chain,
$\cumval > \pay$, creating an asymmetry where the seller's bond exceeds
the buyer's.
\end{definition}

\begin{definition}[Commitment]
\label{def:commitment}
A \emph{commitment} is a tuple $\langle B, S, \tau, \pay, h, \sigma,
\delta \rangle$ where $B$ is the buyer identity, $S$ the seller
identity, $\tau$ the unit of account, $\pay$ the payment, $h$ a
content hash of the underlying agreement, $\sigma$ a cryptographic
salt, and $\delta$ a signature expiry deadline. Both parties sign
the commitment.
\end{definition}

\begin{definition}[Process]
\label{def:process}
A \emph{process} $\Pi$ is a tuple $\langle \text{id}, B, \tau,
\cumval, n, \mathcal{O} \rangle$ where $\text{id}$ is a deterministic
identifier derived from the root commitment, $B$ is the root buyer
(constant across all orders), $\tau$ is the denomination (constant),
$\cumval$ is a \emph{monotonic accumulator} tracking total committed
value, $n$ is the count of active orders, and $\mathcal{O}$ is the
set of active order identifiers.
\end{definition}

\begin{remark}[Monotonicity Assumption]
\label{rem:monotonicity-assumption}
We assume the underlying settlement layer maintains $\cumval$ as a
monotonic accumulator: every successful \texttt{commit} on $\Pi$
increases $\cumval$ by exactly the new order's payment $\pay$, and
no operation decreases $\cumval$. This is an assumption about the
settlement layer's bookkeeping discipline, not a property derived
from the mechanism. The equilibrium analyses that follow apply to
any settlement substrate that satisfies it; an implementation must
enforce the assumption for the analyses to apply. The mechanism
itself is substrate-independent: it can be realised on a
cryptographically settled state machine, on a clearing house, on a
trusted ledger, or on any other settlement layer that maintains
the monotone-accumulator discipline.
\end{remark}

\begin{definition}[Buyer Dominance]
\label{def:buyer-dominance}
A process $\Pi$ exhibits \emph{buyer dominance} when (i) extension
of $\Pi$ via \texttt{commit} requires a dual-signed commitment
authenticated by $B$'s signature, and (ii) only $B$ may trigger
\texttt{resolveProcess}. No other identity---seller, third
party, administrator, or governance body---has authority to
extend or resolve $\Pi$.

The two conditions gate on different identity surfaces. Condition
(i) gates on the identity of the \emph{signer}, not on the agent
that conveys the message: an authorized delegate carrying $B$'s
signature may convey a \texttt{commit} on $B$'s behalf without
breaking buyer dominance, since the substrate verifies the
signature and admits the commitment if and only if it recovers to
$B$. Condition (ii), in contrast, gates on the identity of the
\emph{caller} that issues \texttt{resolveProcess} directly. In
deployed settings $B$ may delegate the resolve action only through
mechanisms that bind delegated authority to $B$'s identity at the
substrate layer (multi-signature wallets, account abstraction, or
key-recovery schemes), not through signature-only conveyance.
\end{definition}

\begin{definition}[Atomic Resolution]
\label{def:atomic-resolution}
Resolution is \emph{atomic}: when the buyer calls
\texttt{resolveProcess}, every active order in the process settles
simultaneously, or none does. There is no partial resolution and
no per-order resolution path. Combined with buyer dominance,
atomic resolution is the second mechanism: it operates on the
mesh of independently bonded edges induced by asymmetric bonding
and makes multi-party coordination resolvable from a single
signature, while inducing the inter-seller weakest-link subgame
analyzed in \Cref{sec:atomic-resolution}.
\end{definition}

\begin{remark}[Single-Currency Binding]
\label{rem:single-currency}
Every order in a process shares a single denomination $\tau$
(\Cref{def:process}). This is not a notational convenience: the
$2\times$ Nash-stability result of \Cref{thm:two-party-nash}
relies on same-unit comparability between the buyer's bond
$\bbond = 2\pay$ and the seller's bond $\sbond = 2\cumval$.
Mixing denominations within one process would require an external
rate of substitution (an oracle, an exchange, or a pre-agreed
conversion rate), each of which introduces a discretionary actor
outside the bond structure and re-opens the escape-hatch problem
of \Cref{thm:escape-hatch-weakness}. Multi-denomination workflows
are accordingly handled outside the mechanism, by parties
converting into the process denomination before issuing a
\texttt{commit}.
\end{remark}

\subsection{Protocol Operations}

The mechanism exposes exactly two state-changing operations.

\paragraph{Operation 1: \texttt{commit}.}
Given a commitment $c$ signed by both $B$ and $S$:
\begin{enumerate}
  \item Verify both signatures.
  \item If $c.\text{processId} = 0$, derive a new process identifier
    from the commitment digest and initialize the process with
    $\cumval \gets \pay$ and $n \gets 1$.
  \item Otherwise verify that $c.\text{processId}$ identifies an
    existing process, that the buyer field $c.B$ matches the process
    root buyer $\Pi.B$ and that the buyer signature recovers to
    $c.B$ (\emph{buyer dominance}), that the denomination
    $c.\tau$ matches $\Pi.\tau$ (\emph{single-currency binding},
    \Cref{rem:single-currency}), and that
    $c.\text{expectedCumulativeValue} = \Pi.\cumval + \pay$
    (\emph{monotonicity}, \Cref{rem:monotonicity-assumption}).
  \item Pull $\bbond = 2\pay$ from $B$ and
    $\sbond = 2c.\text{expectedCumulativeValue}$ from $S$.
  \item Create or extend the process accordingly, update the monotonic
    accumulator, and emit an evidence event that fully records the
    commitment.
\end{enumerate}

\paragraph{Operation 2: \texttt{resolveProcess}.}
Given a process identifier and the complete list of active commitments:
\begin{enumerate}
  \item Verify the caller identity equals $\Pi.B$ (buyer dominance).
  \item Verify that the list length equals $\Pi.n$ (atomic
    resolution: \emph{all} orders must be included).
  \item For each order $o_i$ with payment $\pay_i$ and cumulative
    snapshot $\cumval_i$:
    \begin{align}
      \pi_{s,i} &= 2\cumval_i + \pay_i, \label{eq:seller-payout} \\
      \pi_{b,i} &= \pay_i. \label{eq:buyer-payout}
    \end{align}
  \item Transfer $\pi_{s,i}$ to each seller, $\pi_{b,i}$ to the buyer.
  \item Mark all orders as resolved.
\end{enumerate}

\begin{remark}[Conservation Law]
\label{rem:conservation}
For each order, the total distributed equals total locked:
\begin{equation}
  \pi_{s,i} + \pi_{b,i} = (2\cumval_i + \pay_i) + \pay_i = 2\cumval_i + 2\pay_i = C_{s,i} + C_{b,i}.
\end{equation}
The escrow holds zero balance after resolution.
\end{remark}

\begin{remark}[No Escape Hatches]
\label{rem:no-escape}
There is no timeout, no cancel-after-commitment, no admin override,
no governance vote, and no unilateral exit from the bonded state.
Once both parties have committed, the only path to fund recovery is
buyer-initiated resolution. This is not a temporary simplification; it
is a \emph{security requirement.}
\Cref{thm:escape-hatch-weakness} proves that any additional exit path
weakens the Nash equilibrium.
\end{remark}

\subsection{Net Payoffs}

\begin{definition}[Generalized Payout Functions]
\label{def:payout-functions}
For a bond multiplier $k \geq 1$, an order with payment $\pay > 0$
and cumulative value $\cumval \geq \pay$, define the
\emph{settlement payout functions}:
\begin{align}
  \spay(k, \cumval, \pay) &= k \cdot \cumval + \pay,
    \label{eq:general-seller-payout} \\
  \bpay(k, \pay) &= (k - 1) \cdot \pay.
    \label{eq:general-buyer-payout}
\end{align}
These satisfy the conservation law
$\spay + \bpay = k \cdot \cumval + k \cdot \pay = \sbond + \bbond$
for $\sbond = k \cdot \cumval$ and $\bbond = k \cdot \pay$. The
\emph{net payoff} is the payout minus the deposited bond:
\begin{align}
  \us^{\text{net}}(k) &= \spay - \sbond = \pay,
    \label{eq:seller-net} \\
  \ub^{\text{net}}(k) &= \bpay - \bbond = -\pay.
    \label{eq:buyer-net}
\end{align}
\end{definition}

\begin{remark}
For the mechanism's $k = 2$: $\spay = 2\cumval + \pay$,
$\bpay = \pay$, $\sbond = 2\cumval$, $\bbond = 2\pay$. The
net payoffs ($+\pay$ for the seller, $-\pay$ for the buyer) are
independent of $k$---the multiplier affects only the capital at risk,
not the final value transfer.
\end{remark}

% ════════════════════════════════════════════════════════════════════
\section{Game-Theoretic Analysis}
\label{sec:game-theory}

\subsection{Two-Party Nash Equilibrium}

\begin{definition}[The Bonding Game]
\label{def:bonding-game}
The \emph{bonding game} $\Gamma = (\{B, S\}, \{A_B, A_S\}, \{\ub, \us\})$
is a one-shot, simultaneous-move, complete-information game:
\begin{itemize}
  \item \textbf{Players}: buyer $B$, seller $S$.
  \item \textbf{Action sets}: $A_B = A_S = \{\Coop, \Def\}$.
  \item \textbf{Interpretation}: $\Coop$ for $S$ means delivering the
    agreed goods/services; $\Coop$ for $B$ means calling
    \texttt{resolveProcess}. $\Def$ means failing to perform.
  \item \textbf{Payoff functions} (net of bond deposits), for a root
    order with payment $\pay > 0$ and $\cumval = \pay$:
\end{itemize}
\begin{equation}\label{eq:utility-buyer}
  \ub(a_B, a_S) =
  \begin{cases}
    -\pay & \text{if } a_B = \Coop \text{ and } a_S = \Coop, \\
    -2\pay & \text{otherwise}.
  \end{cases}
\end{equation}
\begin{equation}\label{eq:utility-seller}
  \us(a_B, a_S) =
  \begin{cases}
    +\pay & \text{if } a_B = \Coop \text{ and } a_S = \Coop, \\
    -2\cumval & \text{otherwise}.
  \end{cases}
\end{equation}
The ``otherwise'' payoff is identical across all non-cooperative
outcomes because any defection prevents resolution: bonds remain
locked indefinitely, yielding a net loss equal to the full deposit.
\end{definition}

\begin{table}[ht]
\centering
\caption{Payoff matrix for the bonding game $\Gamma$.
Entry format: $(\ub, \us)$. Parameters: $\pay > 0$,
$\cumval \geq \pay > 0$.}
\label{tab:payoff-matrix}
\begin{tabular}{lcc}
\toprule
 & $a_B = \Coop$ & $a_B = \Def$ \\
\midrule
$a_S = \Coop$ & $(-\pay,\; +\pay)$ & $(-2\pay,\; -2\cumval)$ \\
$a_S = \Def$  & $(-2\pay,\; -2\cumval)$ & $(-2\pay,\; -2\cumval)$ \\
\bottomrule
\end{tabular}
\end{table}

\begin{definition}[Dominance]
\label{def:dominance}
Action $a_i$ \emph{weakly dominates} $a_i'$ for player $i$ if
$u_i(a_i, a_{-i}) \geq u_i(a_i', a_{-i})$ for all $a_{-i}$, with
strict inequality for at least one $a_{-i}$. If the inequality is
strict for \emph{all} $a_{-i}$, the dominance is \emph{strict}.
\end{definition}

\begin{theorem}[Two-Party Nash Equilibrium]
\label{thm:two-party-nash}
In the bonding game $\Gamma$ of \Cref{def:bonding-game}:
\begin{enumerate}[label=(\alph*)]
  \item $\Coop$ weakly dominates $\Def$ for both the buyer and
    the seller.
  \item $(\Coop, \Coop)$ is the unique profile surviving iterated
    elimination of weakly dominated strategies. The profile
    $(\Def, \Def)$ also satisfies the Nash condition (no
    \emph{profitable} deviation, since both yield the locked-bond
    payoff), but is weakly dominated and is removed in the first
    elimination round; we adopt the iterated-elimination-of-weakly-
    dominated-strategies (IEWDS) refinement, standard for
    mechanism design~\citep{myerson1979}, as the operative
    equilibrium concept.
  \item $(\Coop, \Coop)$ is the unique Pareto-optimal outcome.
\end{enumerate}
\end{theorem}

\begin{proof}
\textbf{Part (a): Weak dominance.}

\emph{Buyer.} We compare $\ub(\Coop, a_S)$ vs.\
$\ub(\Def, a_S)$ for each seller action:
\begin{align}
  a_S = \Coop: &\quad \ub(\Coop, \Coop) = -\pay > -2\pay = \ub(\Def, \Coop),
    \label{eq:buyer-strict} \\
  a_S = \Def: &\quad \ub(\Coop, \Def) = -2\pay = \ub(\Def, \Def).
    \label{eq:buyer-weak}
\end{align}
Inequality~\eqref{eq:buyer-strict} is strict (since $\pay > 0$);
\eqref{eq:buyer-weak} is an equality. By \Cref{def:dominance}, $\Coop$
weakly dominates $\Def$ for $B$.

\emph{Seller.} We compare $\us(a_B, \Coop)$ vs.\
$\us(a_B, \Def)$ for each buyer action:
\begin{align}
  a_B = \Coop: &\quad \us(\Coop, \Coop) = +\pay > -2\cumval = \us(\Coop, \Def),
    \label{eq:seller-strict} \\
  a_B = \Def: &\quad \us(\Def, \Coop) = -2\cumval = \us(\Def, \Def).
    \label{eq:seller-weak}
\end{align}
Inequality~\eqref{eq:seller-strict} is strict (since $\pay > 0$ and
$\cumval > 0$ imply $\pay + 2\cumval > 0$);
\eqref{eq:seller-weak} is an equality. $\Coop$ weakly dominates $\Def$
for $S$.

\textbf{Part (b): Unique Nash equilibrium.}
Since $\Coop$ weakly dominates $\Def$ for both players, neither can
profitably deviate from $(\Coop, \Coop)$. We verify no other profile
is a Nash equilibrium:
\begin{itemize}
  \item $(\Coop, \Def)$: Seller deviates to $\Coop$, gaining
    $\pay - (-2\cumval) = \pay + 2\cumval > 0$. Not a NE.
  \item $(\Def, \Coop)$: Buyer deviates to $\Coop$, gaining
    $-\pay - (-2\pay) = \pay > 0$. Not a NE.
  \item $(\Def, \Def)$: Buyer deviates to $\Coop$ with no change
    (both yield $-2\pay$). Seller deviates to $\Coop$ with no change
    (both yield $-2\cumval$). This profile satisfies the Nash
    condition (no \emph{profitable} deviation), but an iterated
    elimination of weakly dominated strategies removes $\Def$ for
    both players, leaving $(\Coop, \Coop)$ as the unique survivor.
\end{itemize}

\textbf{Part (c): Pareto optimality.}
The payoff vector $(-\pay, +\pay)$ Pareto-dominates $(-2\pay, -2\cumval)$
since $-\pay > -2\pay$ and $+\pay > -2\cumval$.
\end{proof}

\begin{remark}[On Weak vs.\ Strict Dominance]
\label{rem:dominance-note}
The dominance is weak, not strict, because when the counterparty
defects, both actions yield the same locked-capital payoff. This is
intrinsic to the mechanism: once one party defects, the outcome is
fully determined regardless of the other's choice. The critical
property is that $(\Def, \Def)$ does not survive iterated elimination
of weakly dominated strategies---the standard refinement for
mechanism design~\citep{myerson1979}.
\end{remark}

\subsection{Sufficiency of the $2\times$ Multiplier}

We now characterize the multiplier $k$ in the generalized bonding game
where $\sbond = k \cdot \cumval$, $\bbond = k \cdot \pay$, and
payouts are given by \Cref{def:payout-functions}.

\begin{lemma}[Buyer Indifference at $k = 1$]
\label{lem:k1-buyer}
For $k = 1$, the buyer is indifferent between cooperating and defecting
when the seller cooperates:
$\ub(\Coop, \Coop) = \ub(\Def, \Coop) = -\pay$.
\end{lemma}
\begin{proof}
With $k = 1$: $\bbond = \pay$ and $\bpay = (k-1)\pay = 0$ (from
\eqref{eq:general-buyer-payout}). Under our accounting convention
(\Cref{def:payout-functions}), the buyer's net payoff under cooperation
is $\ub(\Coop, \Coop) = \bpay - \bbond = 0 - \pay = -\pay$. Under
defection, the buyer's bond remains locked:
$\ub(\Def, \Coop) = -\bbond = -\pay$. Since
$\ub(\Coop, \Coop) = \ub(\Def, \Coop) = -\pay$, the buyer has no
strictly profitable deviation from $\Def$---the mechanism provides
zero marginal incentive to resolve.
\end{proof}

\begin{theorem}[Minimal Viable Bond Multiplier]
\label{thm:2x-sufficient}
Let $k \in \mathbb{R}_{\geq 1}$ be the bond multiplier. Then:
\begin{enumerate}[label=(\alph*)]
  \item For $k = 1$, $\Coop$ does not weakly dominate $\Def$ for the
    buyer: the buyer is indifferent when the seller cooperates.
  \item For every $k > 1$, $\Coop$ weakly dominates $\Def$ for both
    players (\Cref{thm:two-party-nash}); $k = 1$ is therefore the
    infimum of the multipliers for which weak dominance fails on the
    buyer side.
  \item Among integer multipliers, $k = 2$ is the smallest value at
    which weak dominance holds. For any $k > 2$, the equilibrium
    properties are identical to $k = 2$, but capital requirements
    increase by $(k - 2)(\cumval + \pay)$ per order with no
    improvement in dominance relations.
\end{enumerate}
Thus $k = 2$ is the minimal \emph{integer} multiplier that yields
weak dominance of $\Coop$ for both players, and is Pareto-optimal in
capital efficiency among all integer multipliers; the underlying
real-line cutoff is $k = 1$. The integer restriction reflects the
discrete denomination of bonded transfers in deployment,
not a property of the equilibrium analysis.
\end{theorem}

\begin{proof}
\textbf{Part (a).} See \Cref{lem:k1-buyer}. With $k = 1$, the buyer
recovers $\bpay = 0$ upon resolution, having deposited $\bbond = \pay$.
The net monetary payoff from cooperating is $-\pay$; the net from
defecting is also $-\pay$ (bond locked). The buyer cannot be motivated
by the mechanism alone to resolve.

\textbf{Part (b).} For $k > 1$: $\bpay = (k-1)\pay$, $\bbond = k\pay$.
Cooperation yields $\ub(\Coop, \Coop) = (k-1)\pay - k\pay = -\pay$;
defection yields $\ub(\Def, \Coop) = -k\pay$. The cooperation--defection
gap is $-\pay - (-k\pay) = (k-1)\pay$, which is strictly positive for
every $k > 1$. Weak dominance for the buyer therefore holds for every
$k > 1$ and fails at $k = 1$ (the gap collapses to zero), establishing
$k = 1$ as the infimum at which weak dominance fails on the buyer
side. The seller-side weak dominance follows from
\Cref{thm:two-party-nash}(a) for any $k \geq 1$ (the seller's net
payoff under cooperation is $+\pay$ regardless of $k$, while
defection forfeits the bond $k\cumval$).

\textbf{Part (c).} For integer $k$, the smallest value at which
both Part (a) (failure at $k=1$) and Part (b) (success for $k > 1$)
admit a feasible bond is $k = 2$. The surplus capital locked at
$k > 2$, $(k - 2)\pay$ for the buyer and $(k - 2)\cumval$ for the
seller, earns zero return and is pure deadweight cost.
\end{proof}

\subsection{Escape-Hatch Weakness}

\begin{theorem}[Escape-Hatch Weakness]
\label{thm:escape-hatch-weakness}
Let $\Gamma_E$ be the bonding game augmented with a unilateral exit
action $E$ available to player $i$, returning fraction $\alpha \in (0, 1]$
of player $i$'s bond. Then either:
\begin{enumerate}[label=(\alph*)]
  \item $\Coop$ no longer weakly dominates $\Def$ for at least one
    player (the Nash equilibrium of \Cref{thm:two-party-nash}
    may not survive), or
  \item $E$ requires a decision by a third party $J \notin \{B, S\}$
    whose incentives are not constrained by the bond structure.
\end{enumerate}
\end{theorem}

\begin{proof}
We analyze three exhaustive sub-cases of unilateral exit.

\emph{Case 1: Timeout (automatic exit after duration $T$).}
The action space becomes $A_B = \{\Coop, \Def, \text{Wait}\}$ where
$\text{Wait}$ means the buyer does not resolve and collects
$\alpha \cdot \bbond$ after time $T$. The buyer's payoff under
$\text{Wait}$ against a cooperating seller is:
\begin{equation}
  \ub(\text{Wait}, \Coop) = -\bbond + \alpha \cdot \bbond
    = -2\pay(1 - \alpha).
  \label{eq:timeout-payoff}
\end{equation}
Compare with cooperation: $\ub(\Coop, \Coop) = -\pay$. The buyer
prefers $\Coop$ to $\text{Wait}$ iff:
\begin{equation}
  -\pay > -2\pay(1 - \alpha)
  \quad\Longleftrightarrow\quad
  2(1 - \alpha) > 1
  \quad\Longleftrightarrow\quad
  \alpha < \tfrac{1}{2}.
  \label{eq:timeout-threshold}
\end{equation}
For $\alpha \geq \frac{1}{2}$, $\text{Wait}$ yields
$\ub \geq -\pay = \ub(\Coop, \Coop)$, so $\Coop$ no longer
(weakly) dominates $\text{Wait}$. The buyer may rationally wait out the
timeout instead of resolving. Even for $\alpha < \frac{1}{2}$, the
defection cost drops from $2\pay$ to $2\pay(1 - \alpha)$, weakening
the deterrent. The seller, anticipating the timeout, faces a modified
game where the buyer's commitment to resolve is no longer credible---the
timeout provides an alternative path to partial capital recovery,
undermining the ``no escape'' property that sustains the equilibrium.

\emph{Case 2: Arbitrator override.}
A third party $J$ can invoke $E$ on behalf of either player.
Formally, this extends the game to three players $\{B, S, J\}$.
The mechanism's self-enforcing property requires that the payoff
structure \emph{alone} determines the outcome. Introducing $J$
makes the outcome contingent on $J$'s action, and $J$ is not bonded:
$J$ deposits nothing, risks nothing, and has no mechanism-constrained
incentive to decide correctly. Whether $J$ is an individual
(arbitrator), an algorithm (oracle), or a collective (jury), the
resolution path now depends on an agent outside the bond structure.
This is condition (b).

\emph{Case 3: Governance vote.}
A special case of Case~2 where $J$ is replaced by a set of voters
$\{J_1, \ldots, J_m\}$ who collectively decide on $E$. This inherits
all problems of Case~2, plus: (i)~coordination failures among voters,
(ii)~plutocratic capture when voting power is stake-weighted,
(iii)~voter apathy (rational ignorance). Each voter $J_k$ is unbonded.

In all three cases, either $\Coop$ ceases to weakly dominate all
alternatives (Case~1), or resolution depends on an unbonded external
agent (Cases 2--3). \qedhere
\end{proof}

\begin{remark}[On Judicial Review and the Distinction from Cases 2--3]
\label{rem:judicial-review}
The theorem's ``unbonded external agent'' condition targets actors
whose decision is \emph{constitutive} of the resolution path
internal to the mechanism---arbitrators, juries, oracles, voters
embedded in the bonding game itself. It does not target the
external legal forum that may, on the evidence record, adjudicate
a dispute \emph{outside} the bonding game under doctrines of
duress, frustration, impossibility, unconscionability, or public
policy. Those forums are constrained by their own institutional
bond structures (appellate review, professional sanction, doctrinal
consistency) and operate on the bonded commitment as evidentiary
input rather than as a resolution path within the mechanism.
\end{remark}

\begin{corollary}\label{cor:no-timeout}
No timeout duration $T$ and no recovery fraction $\alpha > 0$ can be
added to the bonding game without weakening the buyer's incentive to
cooperate. In particular, full recovery ($\alpha = 1$) reduces the
game to a zero-cost option, making defection strictly dominant for
the buyer.
\end{corollary}

\begin{remark}[Robustness to Weaker Rationality]
\label{rem:bounded-rationality}
The dominance results stated in monetary terms extend to a wider
class of agent preferences than risk-neutral expected-payoff
maximisation. Three observations.
First, the dominance argument compares the payoff under cooperation
($+\pay_i$ for $S_i$, $-\pay_i$ for $B$) to the payoff under
defection ($-2\cumval_i$ for $S_i$, $-2\pay_i$ for $B$); since
$\pay_i > -2\cumval_i$ and $-\pay_i > -2\pay_i$ are preserved under
any strictly monotone utility transformation, weak dominance of
$\Coop$ holds for any preference order that prefers more money to
less, including arbitrary risk-averse and loss-averse utility
specifications.
Second, the cooperation--defection gap $\pay_i + 2\cumval_i$ is
bounded away from zero (and grows with chain depth), so the
dominance margin is not perturbation-sensitive: a trembling-hand
refinement preserves $(\Coop, \ldots, \Coop)$ as the unique
surviving profile, since small perturbations to other players'
actions cannot flip a strictly positive gap.
Third, the dominance argument requires only one comparison per
player and does not depend on iterated reasoning about counterparty
rationality, common knowledge of payoffs, or beliefs about types.
What remains genuinely outside the analysis is behaviour under
non-pecuniary preferences---spite, fairness norms, or intrinsic
motivation to defect---where the monetary payoff matrix does not
capture all relevant utilities.
\end{remark}

% ════════════════════════════════════════════════════════════════════
\section{$N$-Party Process Chains}
\label{sec:process-trees}

\subsection{$N$-Party Asymmetric Bonding}

We now extend the two-party model to an arbitrary-length chain of
sellers.

\begin{definition}[Process Chain]
\label{def:process-chain}
A \emph{process chain} is a sequence of bonded orders
$\langle o_1, o_2, \ldots, o_n \rangle$ sharing a common buyer $B$ and
denomination $\tau$. Order $o_i$ has local payment $\pay_i > 0$ and
cumulative value:
\begin{equation}
  \cumval_i = \sum_{j=1}^{i} \pay_j.
  \label{eq:cumulative}
\end{equation}
The seller of $o_i$ posts bond $C_{s,i} = 2\cumval_i$; the buyer
posts bond $C_{b,i} = 2\pay_i$.
\end{definition}

\begin{definition}[$N$-Party Bonding Game]
\label{def:n-party-game}
The \emph{$N$-party bonding game}
$\Gamma_N = (\{B, S_1, \ldots, S_n\}, \{A_i\}, \{u_i\})$
extends $\Gamma$ to $n$ sellers. Each seller $S_i$ has action set
$\{\Coop, \Def\}$. The buyer has action set $\{\Coop, \Def\}$ and
chooses after observing each seller's realized action; we assume
perfect monitoring, so $B$ perfectly observes whether each $S_i$
delivered, and plays the conditional strategy ``resolve iff all
sellers cooperated.'' The perfect-monitoring assumption
corresponds to the deployed setting, in which the buyer receives
either the agreed goods or not, and the buyer's observation of
delivery is prior to the buyer's \texttt{resolveProcess} call.
Imperfect monitoring is addressed at composable surfaces above
the mechanism (attestation, schema-typed witness records,
third-party signal authorization), which lie outside the present
paper's scope. Payoffs under perfect monitoring:
\begin{equation}\label{eq:n-party-seller-utility}
  u_{S_i}(a_B, a_{S_1}, \ldots, a_{S_n}) =
  \begin{cases}
    +\pay_i & \text{if } a_B = \Coop \text{ and } a_{S_j} = \Coop
      \;\forall j, \\
    -2\cumval_i & \text{otherwise}.
  \end{cases}
\end{equation}
The buyer's payoff is:
\begin{equation}\label{eq:n-party-buyer-utility}
  \ub(a_B, a_{S_1}, \ldots, a_{S_n}) =
  \begin{cases}
    -\sum_{i=1}^{n} \pay_i & \text{if } a_B = \Coop \text{ and }
      a_{S_j} = \Coop \;\forall j, \\
    -\sum_{i=1}^{n} 2\pay_i & \text{otherwise}.
  \end{cases}
\end{equation}
\end{definition}

\begin{remark}[Sequential vs Simultaneous Equivalence]
\label{rem:sequential-simultaneous}
The game $\Gamma_N$ is presented in simultaneous-move form, but in
deployment the events of \texttt{commit}, delivery, and
\texttt{resolveProcess} are sequential. The two representations are
equivalent under asymmetric bonding: each seller $S_i$'s payoff
function depends on $S_i$'s own action and the buyer's eventual
\texttt{resolve}/\texttt{not-resolve} decision, and not on the
order in which other sellers acted or on the temporal position of
$S_i$'s \texttt{commit} relative to others'. Defection by any
participant forces non-resolution
(\Cref{def:atomic-resolution}), which yields the same locked-bond
payoff $-2\cumval_i$ to $S_i$ regardless of when in the sequence
the defection occurs. The dominance argument of
\Cref{thm:two-party-nash} therefore applies edge-by-edge at every
node of the corresponding extensive-form game, so the IEWDS
profile of the simultaneous-move representation coincides with
the subgame-perfect Nash equilibrium of the extensive-form game.
We use the simultaneous-move representation for expositional
clarity.
\end{remark}

\begin{lemma}[Buyer Best Response]
\label{lem:buyer-best-response}
The buyer's \emph{post-commit} action set in $\Gamma_N$ is
$A_B = \{\text{resolve}, \text{not-resolve}\}$ (the buyer's
prior \texttt{commit} actions establish the chain and are
fixed at the point this lemma applies). Under perfect monitoring,
the buyer's best response to the cooperative seller profile
$(\Coop, \ldots, \Coop)$ is $\text{resolve}$. Under any profile
in which at least one seller defects, the buyer is indifferent
between $\text{resolve}$ and $\text{not-resolve}$ (both yield
the locked-bond payoff $-\sum_i 2\pay_i$).
\end{lemma}

\begin{proof}
The buyer's only post-commit decision is whether to call
\texttt{resolveProcess}; there is no continuation action that
strictly improves on $-\sum_i 2\pay_i$ once any seller has
defected, by atomic resolution
(\Cref{def:atomic-resolution}). Under $(\Coop, \ldots, \Coop)$,
the buyer's payoff under $\text{resolve}$ is $-\sum_i \pay_i$
versus $-\sum_i 2\pay_i$ under $\text{not-resolve}$; the
cooperation--defection gap $\sum_i \pay_i > 0$ is the linear
sum of the bilateral gaps of \Cref{thm:two-party-nash}(a),
applied to each order independently. Hence $\text{resolve}$
strictly dominates $\text{not-resolve}$ on the cooperative
profile, and the two are payoff-equal off it.
\end{proof}

\begin{theorem}[$N$-Party Nash Equilibrium]
\label{thm:n-party-nash}
In $\Gamma_N$, cooperation is a weakly dominant strategy for every
seller $S_i$ at every position $i \in \{1, \ldots, n\}$, and
$(\Coop, \Coop, \ldots, \Coop)$ is the unique outcome surviving
iterated elimination of weakly dominated strategies.
\end{theorem}

\begin{proof}
Fix an arbitrary seller $S_i$. Consider the two cases for the remaining
players' joint action profile $a_{-i}$.

\emph{Case 1: All other sellers cooperate and buyer resolves.} Then:
\begin{align}
  u_{S_i}(\ldots, \Coop_i, \ldots) &= +\pay_i, \\
  u_{S_i}(\ldots, \Def_i, \ldots) &= -2\cumval_i.
\end{align}
Since $\pay_i > 0$ and $\cumval_i > 0$, we have
$\pay_i > -2\cumval_i$, i.e., $\pay_i + 2\cumval_i > 0$. The
cooperation payoff strictly exceeds the defection payoff.

\emph{Case 2: At least one other seller defects, or buyer defects.}
Then resolution does not occur regardless of $S_i$'s action:
\begin{equation}
  u_{S_i}(\ldots, \Coop_i, \ldots) = u_{S_i}(\ldots, \Def_i, \ldots)
    = -2\cumval_i.
\end{equation}
Payoffs are equal.

Combining both cases, $\Coop$ weakly dominates $\Def$ for $S_i$, with
strict inequality in Case~1. Since the choice of $i$ was arbitrary,
this holds for all $n$ sellers. The buyer's dominance follows from
\Cref{lem:buyer-best-response}: on the cooperative seller profile
the buyer's best response is $\text{resolve}$, with payoff
$-\sum \pay_i$ strictly exceeding the payoff $-2\sum \pay_i$ under
$\text{not-resolve}$.

Iterated elimination: in the first round, $\Def$ is eliminated for all
players (weakly dominated). The unique surviving profile is
$(\Coop, \ldots, \Coop)$.
\end{proof}

\begin{remark}[Strength of the $N$-Party Result]
The seller's dominance condition,
$\pay_i + 2\cumval_i > 0$, holds with a margin that \emph{grows}
with chain depth (since $\cumval_i = \sum_{j \leq i} \pay_j$ is
monotonic). The cooperation-defection payoff gap for $S_i$ is:
\begin{equation}
  \Delta_i = \pay_i - (-2\cumval_i) = \pay_i + 2\cumval_i
    = \pay_i + 2\!\sum_{j=1}^{i} \pay_j \geq 3\pay_i.
\end{equation}
The gap is at least $3\pay_i$ (equality only at the root where
$\cumval_i = \pay_i$) and grows with every upstream payment. This
is the formal basis of the ``geometric coordination pressure'' in
\Cref{cor:geometric-pressure}.
\end{remark}

\begin{corollary}[Coordination Pressure Grows With Depth]
\label{cor:geometric-pressure}
The risk-to-reward ratio for seller $S_i$ is:
\begin{equation}
  \rho_i = \frac{C_{s,i}}{\pay_i} = \frac{2\cumval_i}{\pay_i}
    = \frac{2\sum_{j=1}^{i} \pay_j}{\pay_i}.
\end{equation}
For equal payments ($\pay_j = p$ for all $j$), $\rho_i = 2i$, growing
linearly with chain depth. For value chains where early stages
contribute disproportionately ($\pay_1 \gg \pay_n$, with $\pay_n$
bounded below), $\rho_n$ scales as $2 \cumval_n / \pay_n$ and grows
faster than linearly in $n$ relative to the late-stage contribution.
The qualitative point is that the deeper participant absorbs the
cumulative upstream risk; the precise growth rate depends on the
shape of the payment schedule.
\end{corollary}

\begin{example}
\label{ex:food-delivery}
Alice orders food (\$10) from Bob, who needs delivery (\$2) from
Charlie:
\begin{center}
\begin{tabular}{lccccc}
\toprule
Party & Role & $\pay_i$ & $\cumval_i$ & $C_{s,i}$ & $\rho_i$ \\
\midrule
Bob     & Seller 1 & \$10 & \$10 & \$20 & 2.0 \\
Charlie & Seller 2 & \$2  & \$12 & \$24 & 12.0 \\
\bottomrule
\end{tabular}
\end{center}
Charlie risks \$24 to earn \$2 ($\rho = 12$). Bob risks \$20 to earn
\$10 ($\rho = 2$). The deeper participant faces amplified consequences.
\end{example}

\subsection{Self-Enforcing Cumulative-Value Reporting}
\label{sec:self-enforcing-reporting}

A natural question: must the mechanism validate the cumulative value
reported by each seller at commit time? It need not. The bonding
rule itself makes honest reporting the dominant strategy whenever
the probability of non-resolution is positive.

\begin{theorem}[Cumulative-Value Reporting Honesty]
\label{thm:self-enforcing-reporting}
In any game where seller $S_i$ chooses a reported cumulative value
$\cumval'_i \geq \cumval_i$ (with $\cumval_i$ being the true value
maintained by the settlement layer's accumulator under
\Cref{rem:monotonicity-assumption}), truth-telling
($\cumval'_i = \cumval_i$) weakly dominates all alternatives. The
dominance is strict whenever the probability of non-resolution is
positive.
\end{theorem}

\begin{proof}
By \Cref{rem:monotonicity-assumption} the accumulator constraint
$\cumval'_i \geq \cumval_i \coloneqq \Pi.\cumval_{\text{prev}} + \pay_i$
holds: any report with $\cumval'_i < \cumval_i$ is rejected by the
settlement layer at \texttt{commit} time and never enters the
accumulator (cf.\ \Cref{def:process,def:commitment}, the
\texttt{expectedCumulativeValue} clause). We therefore restrict
attention to $\cumval'_i \geq \cumval_i$.

Let $q \in [0, 1]$ be the probability that the buyer resolves the
process (a function of all sellers' performance, exogenous to $S_i$'s
reporting decision). Define $S_i$'s expected payoff as a function of
the report $\cumval'_i$:
\begin{equation}
  \mathbb{E}[u_{S_i}(\cumval'_i)]
    = q \cdot (\underbrace{2\cumval'_i + \pay_i}_{\text{payout}}
             - \underbrace{2\cumval'_i}_{\text{bond}})
    + (1 - q) \cdot (-2\cumval'_i)
    = q \cdot \pay_i - (1 - q) \cdot 2\cumval'_i.
  \label{eq:reporting-expected-utility}
\end{equation}
The derivative with respect to $\cumval'_i$ is:
\begin{equation}
  \frac{\partial}{\partial \cumval'_i}
    \mathbb{E}[u_{S_i}] = -2(1 - q).
  \label{eq:reporting-derivative}
\end{equation}
For $q < 1$ (non-zero probability of non-resolution), this derivative
is strictly negative: expected utility is strictly decreasing in the
reported cumulative value. The optimum is therefore at the smallest
feasible value, $\cumval'_i = \cumval_i$.

For $q = 1$ (certain resolution), the derivative is zero and the
seller is indifferent---overstating wastes capital but does not
affect the net payoff. Truth-telling remains weakly optimal.

Since $q < 1$ in any realistic setting (the seller cannot guarantee
other sellers' cooperation), truth-telling strictly dominates
overstating.
\end{proof}

\begin{remark}[Why the Mechanism Does Not Validate Cumulative-Value Reports]
The theorem licenses a deliberate omission in the mechanism: beyond
enforcing monotonicity of the accumulator, the mechanism performs no
validation of the cumulative-value reports it receives. Misreporting
cannot profit a seller because the bonding rule already makes
overstating $\cumval_i$ strictly dominated.
\end{remark}

\subsection{Atomic Resolution and Seller Coordination}
\label{sec:atomic-resolution}

Atomic resolution was defined formally in \Cref{def:atomic-resolution};
we now develop its game-theoretic consequence.

\begin{theorem}[Endogenous Coordination Pressure]
\label{thm:coordination-pressure}
The $N$-party bonding game under atomic resolution induces a
\emph{weakest-link public goods game} among sellers, where each
seller's payout depends on universal cooperation.
\end{theorem}

\begin{proof}
Consider the sellers' subgame (fixing the buyer's strategy at: resolve
iff all sellers cooperate). Each $S_i$ has payoff
$u_{S_i}$ as in \eqref{eq:n-party-seller-utility}.

Define the \emph{externality} imposed by $S_k$'s defection on $S_i$
($i \neq k$):
\begin{equation}
  \varepsilon_{k \to i}
    = u_{S_i}(\text{all}\;\Coop) - u_{S_i}(S_k\;\text{defects})
    = \pay_i - (-2\cumval_i)
    = \pay_i + 2\cumval_i.
\end{equation}
This is strictly positive for all $i, k$ with $i \neq k$. Every
defection imposes a loss on every other seller. The total externality
of $S_k$'s defection is:
\begin{equation}
  \mathcal{E}_k = \sum_{i \neq k} \varepsilon_{k \to i}
    = \sum_{i \neq k} (\pay_i + 2\cumval_i).
\end{equation}

This structure satisfies the defining properties of a \emph{weakest-link
game}~\citep{hirshleifer1983}:
\begin{enumerate}
  \item The group outcome depends on the minimum individual effort
    (one defection collapses all payoffs to the non-cooperative level).
  \item The cooperative outcome $(\Coop, \ldots, \Coop)$ is Pareto
    optimal among the subgame's Nash equilibria; the all-defect
    profile satisfies the Nash condition (no profitable deviation
    from a locked-bond payoff) but is removed by IEWDS, leaving
    $(\Coop, \ldots, \Coop)$ as the unique IEWDS-surviving profile.
  \item The action structure is \emph{weak-link best-response
    coupling under perfect monitoring}: $S_i$'s cooperation--defection
    gap is $\pay_i + 2\cumval_i$ when every other seller and the
    buyer cooperate, and zero otherwise. The marginal value of
    cooperation is therefore not monotonically increasing in others'
    cooperation (the standard strategic-complementarity shape); it
    is constant on the all-cooperate profile and zero elsewhere,
    which is the characteristic structure of weak-link coordination
    rather than strategic complementarity.
\end{enumerate}

Under one-shot, anonymous interaction, this externality structure
suffices for the cooperation-pressure result of
\Cref{prop:microfinance-reduction}: every $S_i$ has a direct
mechanism-level economic interest of magnitude $\pay_i + 2\cumval_i$
in $S_k$'s cooperation, computable from the mechanism's monotonic
accumulator alone. No information acquisition, repeated interaction,
or social-substrate communication is required for the externality
to be present; whether parties additionally engage in such
practices in deployment is orthogonal to the equilibrium argument.
\end{proof}

\begin{remark}[Distinction from the Experimental Weakest-Link Literature]
\label{rem:vanhuyck}
The experimental weakest-link literature, originating with
\citet{vanhuyck1990},
finds that weakest-link games empirically produce coordination
\emph{failure} in groups larger than two: players converge to the
Pareto-inferior all-defect equilibrium via pessimistic beliefs about
co-players. This might appear to threaten the result above. It does
not. The classical experimental result arises when cooperation is
\emph{collectively} preferred but not \emph{individually} dominant ---
a setting where pessimistic beliefs about others are sufficient to
rationalise defection. In the bonded commitment mechanism, defection
is individually dominated regardless of what others do (Case~2 of the
proof of \Cref{thm:n-party-nash}: even when all others defect, $S_i$
incurs $-2\cumval_i$ whether it cooperates or defects, so there is no
strategic incentive to defect). The cooperation pressure identified in
\Cref{thm:coordination-pressure} therefore arises in a setting where
coordination failure through pessimistic beliefs cannot occur: each
seller's dominant strategy is already to cooperate, and the weakest-link
structure adds the externality characterisation on top of that result,
not as a substitute for it.
\end{remark}

\begin{proposition}[Collusion is Strictly Unprofitable]
\label{prop:collusion-unprofitable}
Let $K \subseteq \{S_1, \ldots, S_n\}$ be any non-empty coalition
of sellers, and let $X \notin \{B, S_1, \ldots, S_n\}$ be any
external party who pays a side transfer $T \geq 0$ to $K$ in
exchange for joint defection by every member of $K$. Joint
defection by $K$ is not a Pareto improvement over universal
cooperation for the coalition $K \cup \{X\}$: the minimum side
transfer at which $K$'s members are individually willing to
defect strictly exceeds the maximum monetary benefit $X$ can
obtain from $B$'s non-resolution. The dominance margin is at
least $3 \sum_{i \in K} \pay_i$ versus a benefit upper-bounded
by $3 \sum_{i \in K} \pay_i$ at the root order, with the gap
strictly positive whenever $K$ contains any non-root seller.
\end{proposition}

\begin{proof}
For each $S_i \in K$, the cooperation--defection gap is
$\Delta_i = \pay_i - (-2\cumval_i) = \pay_i + 2\cumval_i$
(\Cref{thm:n-party-nash}, Case~1). For $S_i$ to weakly accept
the side payment, $X$ must transfer at least $\Delta_i$ to $S_i$.
The minimum total side transfer is therefore
\begin{equation}
  T_{\min}(K) = \sum_{i \in K} \Delta_i
    = \sum_{i \in K} (\pay_i + 2\cumval_i)
    \geq 3 \sum_{i \in K} \pay_i,
\end{equation}
with equality only when every $S_i \in K$ is a root order
($\cumval_i = \pay_i$).

$X$'s monetary benefit from $K$'s joint defection is bounded by
the most an arbitrary external party with spite preferences could
rationally pay for the buyer's harm. That harm has two components.
Non-resolution locks $B$'s bonds, totalling $\sum_{i=1}^n 2\pay_i$,
and denies $B$ the goods themselves; the value of the goods to
$B$ is bounded above by $B$'s willingness to pay,
$\sum_{i=1}^n \pay_i$ (the participation constraint that says
$B$ entered the bonding game). Restricting to the colluded
orders, $X$'s upper bound on benefit is
\begin{equation}
  B_{\max}(K) \leq \sum_{i \in K} 2\pay_i
    + \sum_{i \in K} \pay_i
    = 3 \sum_{i \in K} \pay_i.
\end{equation}

Combining the bounds,
$T_{\min}(K) \geq 3\sum_{i \in K} \pay_i \geq B_{\max}(K)$,
with strict inequality on the first step whenever $K$ contains
a seller $S_i$ with $\cumval_i > \pay_i$ (i.e., any non-root
seller). For collusion to be Pareto-improving for $K \cup \{X\}$,
$X$ would need benefit strictly greater than the side-payment
cost; this is impossible. Joint defection is therefore strictly
Pareto-dominated.
\end{proof}

\begin{remark}[Margin Grows with Chain Depth]
\label{rem:collusion-depth}
The dominance margin $T_{\min}(K) - B_{\max}(K) \geq
2 \sum_{i \in K} (\cumval_i - \pay_i) \geq 0$ is monotone in
chain depth: the deeper the colluded sellers sit in the chain,
the larger the spread between their bonds $2\cumval_i$ and the
harm imposable on the buyer. Coalitions of deep-chain sellers
are accordingly the most strictly unprofitable.
\end{remark}

\subsection{Seller Non-Performance and Externality Handling}
\label{sec:externalities}

A natural question in any multi-party coordination mechanism is how
externalities are handled. Four categories arise, and the bonded
commitment mechanism treats them with a single primitive ---
bonded sub-orders, optionally layered with external risk products
--- without requiring additional governance or protocol
modifications.

\begin{enumerate}[label=(\roman*)]
  \item \emph{Intra-process externalities}---the loss one seller's
    defection imposes on co-sellers---are handled by the bonding
    mechanism itself: defection is weakly dominated for every seller
    at every position (\Cref{thm:n-party-nash}), and the cooperation
    pressure that propagates through atomic resolution
    (\Cref{thm:coordination-pressure}) is the externality made
    strategic. No additional handling is required.
  \item \emph{Seller non-performance}---the focus of the two
    paragraphs that follow---is absorbed either by the buyer
    committing a substitute order before resolution, or by an
    external risk product (performance guarantee, surety bond,
    insurance) layered over the bonded position.
  \item \emph{Buyer non-performance}---the buyer commits but
    never resolves---is absorbed by the no-escape-hatch property:
    the buyer's bond remains locked indefinitely, imposing a
    $2\pay$ deterrent equal to the buyer's worst-case loss
    (\Cref{rem:no-escape}, \Cref{thm:escape-hatch-weakness}). The
    sellers' bonds are likewise locked indefinitely; recovery for
    sellers in this case requires the same external risk products
    as in (ii).
  \item \emph{Third-party externalities}---harms imposed by the
    bonded process on non-participants (a courier causing a
    traffic accident, a kitchen polluting, a process generating
    a data leak)---compose into the mechanism the same way as
    (ii). The harmed party becomes a remediation buyer in a new
    bonded sub-order against the harming party as remediation
    seller; or an external risk product (liability insurance,
    pollution bond) layers over the bonded process in advance.
    The mechanism is uniform across externality types: the same
    composition pattern handles internal and third-party
    externalities, with the harmed party simply being internal
    or external to the original process.
\end{enumerate}

The two paragraphs below develop case (ii) in detail; cases (i),
(iii), and (iv) are addressed by cross-reference and by the same
composition pattern.

\paragraph{Seller substitution.}
If seller $S_i$ fails to perform, the buyer holds a non-resolved
process in which $S_i$'s bond is locked but delivery has not
occurred. Before calling \texttt{resolveProcess}, the buyer may
commit a replacement order with a new seller $S_i'$ at the same or
renegotiated payment. $S_i'$ bonds against the current cumulative
value, enters the process under the standard bonding rule, and
delivers. When the buyer resolves, $S_i'$ receives full settlement.
$S_i$, having failed to perform, forfeits their bond under the
bonding equilibrium. The externality of $S_i$'s non-performance is
therefore absorbed by $S_i$'s own capital loss, not distributed to
co-sellers: other sellers' payouts are unaffected.

\paragraph{External risk products.}
Where seller substitution is impractical---latency-sensitive
processes, highly specialised roles, or cases where a replacement
cannot be sourced---standard risk instruments such as performance
guarantees, surety bonds, and insurance can layer over the bonded
process without modifying the mechanism. These instruments hold
independently bonded positions on the seller's behalf; they do not
call \texttt{commit} or \texttt{resolveProcess}, hold no mechanism
state, and leave the bonding equilibrium intact. The mechanism is
therefore not the only layer of externality protection available to
parties; it is the settlement primitive over which such risk
products compose.

\subsection{Comparison with Joint-Liability Microfinance}
\label{sec:microfinance}

The coordination pressure of \Cref{thm:coordination-pressure} is
operationally analogous to the peer-enforcement equilibrium documented
in joint-liability microfinance, formalized by \citet{ghatak1999} and
\citet{stiglitz1990} on the Grameen Bank model~\citep{yunus1999}.
Those models obtain peer enforcement through a combination of
joint-liability contracting and a social substrate (repeated
interaction, local information, exogenous punishment via social
sanction); the contract alone is insufficient, and the substrate
does load-bearing enforcement work. The bonded commitment mechanism
reproduces the same equilibrium-selection effect under strictly
weaker assumptions.

\begin{proposition}[Assumption Reduction on Cooperation Pressure]
\label{prop:microfinance-reduction}
The \emph{cooperation-pressure component} of the equilibrium common
to \citet{ghatak1999} and \citet{stiglitz1990}---cooperative
equilibrium selection driven by co-participant externalities---is
obtained by the bonded commitment mechanism under strictly weaker
assumptions. Specifically, \figaro{} requires none of the following:
(i)~repeated interaction,
(ii)~local information among sellers,
(iii)~a punishment technology exogenous to the contract, nor
(iv)~joint-liability contracting.
\end{proposition}

\begin{proof}
\Cref{thm:n-party-nash} establishes cooperation as the unique profile
surviving iterated elimination of weakly dominated strategies in a
one-shot game. The argument uses only the mechanism's payoff function
and makes no appeal to continuation value, trigger strategies, or
folk-theorem constructions; (i) is not required.

The externality $\varepsilon_{k \to i} = \pay_i + 2\cumval_i$
(\Cref{thm:coordination-pressure}) is computed from quantities
enforced by the settlement layer's monotonic accumulator, independent
of any seller's information about co-sellers' types, actions, or
identities; (ii) is not required.

The punishment for a co-seller's defection is the loss of $S_i$'s own
bond $C_{s,i} = 2\cumval_i$, imposed directly by the mechanism through
non-resolution. No external punishment technology is invoked;
(iii) is not required.

Finally, each $C_{s,i}$ is posted individually by $S_i$ against
$S_i$'s own cumulative-value snapshot, not against the group's
aggregate obligation. Coupling across sellers arises from the buyer's
atomic-resolution constraint (\Cref{def:atomic-resolution}), not from
the bond structure; (iv) is not required.
\end{proof}

The reduction is scoped to cooperation pressure. Two further
microfinance results---peer selection on private types
\citep{ghatak1999} and peer monitoring under principal information
asymmetry \citep{stiglitz1990}---are not directly reproduced; both
require additional structure above the bonded primitive (attested
credentials, signal-contingent resolution rules) and lie outside
the present paper's scope.

The reduction reframes rather than corrects the microfinance theorems:
they remain sound under their assumptions, but the assumptions
characterize one particular settlement environment in which the
mechanism was historically deployed (reputation within a social
network), not necessary conditions for endogenous peer monitoring.
Once the settlement layer admits bilaterally signed, individually
collateralized commitments with atomic resolution, the same outcome
obtains among mutual strangers in a single shot with no social
substrate.

% ════════════════════════════════════════════════════════════════════
\section{Dispute Resolution}
\label{sec:dispute}

Dispute resolution is marginal in this mechanism by design.
Defection is weakly dominated for every party at every position
(\Cref{thm:two-party-nash}, \Cref{thm:n-party-nash}); the
cooperation pressure that propagates through atomic resolution
(\Cref{thm:coordination-pressure}) sharpens this further by
making each seller's payout contingent on universal cooperation
among co-sellers; and the reduction over joint-liability
microfinance (\Cref{prop:microfinance-reduction}) places the
bonded mechanism in the same equilibrium-selection regime as
group lending under strictly weaker assumptions. The bonded
process itself absorbs almost all of the work that traditional
arrangements externalise to courts, arbitrators, and platform
adjudication.

The microfinance literature provides a qualitative empirical
reference. Joint-liability group lending, beginning with the
Grameen Bank, has produced low documented default rates under
conditions (repeated interaction, local information, social
sanction) that the bonded mechanism replaces with one-shot
collateral and atomic resolution~\citep{yunus1999, ghatak1999}.
The fraction of bonded processes that require external dispute
resolution is therefore expected to be small, but the present
mechanism does not predict the rate quantitatively: the
microfinance default rate emerges from the combined effect of
peer selection, peer monitoring, and cooperation pressure, only
the last of which the bonded mechanism reproduces
(\Cref{prop:microfinance-reduction}). Calibration is an
empirical matter for deployment data, not a theorem of the
present paper.

For the marginal cases where internal pressure does not resolve a
dispute, the bonded commitment is compatible with external legal
and arbitral forums. \Cref{rem:judicial-review} establishes the
formal distinction: such forums adjudicate \emph{outside} the
bonding game on the evidence record produced by the mechanism, rather
than serving as a constitutive resolution path within the
mechanism. Parties may name a forum-clause in their assembly graph
(\Cref{sec:architectural-posture}) selecting Kleros, an arbitral
body, a court, or any other forum~\citep{kleros2019}; the mechanism
does not bind a forum, and the choice does not affect the bonding
equilibrium. The composition is asymmetric in an important way:
\figaro{} composes with external forums as evidentiary consumers
of the bonded record, but external forums cannot compose with
\figaro{} as a constitutive resolution path within the mechanism
without breaking the equilibrium (\Cref{thm:escape-hatch-weakness},
Case~2). The asymmetry is the formal expression of the design
stance: dispute resolution is a marginal external concern, not an
internal mechanism component.

% ════════════════════════════════════════════════════════════════════
\section{Conclusion}
\label{sec:conclusion}

We have presented \figaro{}, a settlement primitive that achieves
self-enforcing multi-party coordination through two composing
mechanisms: \emph{asymmetric bonding}, which produces the bilateral
Nash equilibrium and scales it naturally to $N$-party process chains
by tying each seller's bond to cumulative upstream value; and
\emph{buyer dominance with atomic resolution}, which enforces
inter-seller coordination on the already-scaled mesh. The mechanism is
minimal---no timeout, no governance, no escape hatches---and the
game-theoretic properties are strong: under perfect monitoring
(\Cref{def:n-party-game}) and complete information, cooperation is
weakly dominant at every position in an $N$-party process chain, and
the cooperative profile is the unique outcome surviving iterated
elimination of weakly dominated strategies. The mechanism reproduces
the cooperation-pressure component of the peer-enforcement equilibrium
in joint-liability microfinance under strictly weaker assumptions,
with no appeal to repeated interaction, local information, or social
sanction.

\section*{Acknowledgments}
I thank the Solidity team for the Safe Remote Purchase pattern,
Fabrizio Romano Genovese for his contributions to \figaro{}, and
the broader Ethereum community for the body of work in mechanism
design, coordination economics, and applied cryptographic
settlement that this paper builds on.


% ════════════════════════════════════════════════════════════════════
% References
% ════════════════════════════════════════════════════════════════════

\begin{thebibliography}{99}

\bibitem[Van~Huyck et~al.(1990)]{vanhuyck1990}
Van Huyck, J.~B., Battalio, R.~C., and Beil, R.~O.
\newblock Tacit coordination games, strategic uncertainty, and coordination
  failure.
\newblock \emph{American Economic Review}, 80(1):234--248, 1990.

\bibitem[Aragon(2017)]{aragon2017}
Cuende, L. and Izquierdo, J.
\newblock Aragon Network: A Decentralized Infrastructure for Value Exchange.
\newblock Aragon White Paper, 2017.
\newblock \url{https://github.com/aragon/whitepaper}

\bibitem[Asgaonkar and Krishnamachari(2019)]{asgaonkar2019}
Asgaonkar, A. and Krishnamachari, B.
\newblock Solving the Buyer and Seller's Dilemma: A Dual-Deposit Escrow Smart
  Contract for Provably Cheat-Proof Delivery and Payment for a Digital Good
  without a Trusted Mediator.
\newblock In \emph{Proc.\ IEEE International Conference on Blockchain and
  Cryptocurrency}, 2019.

\bibitem[Zimbeck(2014)]{bithalo2014}
Zimbeck, D.
\newblock Two Party double deposit trustless escrow in cryptographic
  networks and Bitcoin.
\newblock BitHalo White Paper, 2014.

\bibitem[Solidity Team(2016)]{solidity2016}
Solidity Team.
\newblock Safe Remote Purchase.
\newblock In \emph{Solidity Documentation---Solidity by Example}, 2016.
\newblock \url{https://docs.soliditylang.org/en/latest/solidity-by-example.html#safe-remote-purchase}

\bibitem[Ghatak(1999)]{ghatak1999}
Ghatak, M.
\newblock Group lending, local information and peer selection.
\newblock \emph{Journal of Development Economics}, 60(1):27--50, 1999.

\bibitem[Yunus(1999)]{yunus1999}
Yunus, M.
\newblock Banker to the Poor: Micro-Lending and the Battle Against World
  Poverty.
\newblock PublicAffairs, New York, 1999.

\bibitem[Hirshleifer(1983)]{hirshleifer1983}
Hirshleifer, J.
\newblock From weakest-link to best-shot: The voluntary provision of public
  goods.
\newblock \emph{Public Choice}, 41(3):371--386, 1983.

\bibitem[Lesaege et~al.(2019)]{kleros2019}
Lesaege, C., George, W., and Ast, F.
\newblock Kleros: Long Paper v1.0.7.
\newblock Kleros, 2019.
\newblock \url{https://kleros.io/yellowpaper.pdf}

\bibitem[Myerson(1979)]{myerson1979}
Myerson, R.~B.
\newblock Incentive Compatibility and the Bargaining Problem.
\newblock \emph{Econometrica}, 47(1):61--73, 1979.

\bibitem[Optimism(2021)]{optimism2021}
Optimism Foundation.
\newblock Optimistic Rollups.
\newblock Optimism Technical Documentation, 2021.
\newblock \url{https://community.optimism.io/docs/protocol/}

\bibitem[Poon and Dryja(2016)]{poon2016}
Poon, J. and Dryja, T.
\newblock The Bitcoin Lightning Network: Scalable Off-Chain Instant Payments.
\newblock Technical report, 2016.

\bibitem[Raiden(2017)]{raiden2017}
Raiden Network.
\newblock Raiden Network: Fast, Cheap, Scalable Token Transfers for Ethereum.
\newblock Raiden White Paper, 2017.
\newblock \url{https://raiden.network/}

\bibitem[Stiglitz(1990)]{stiglitz1990}
Stiglitz, J.~E.
\newblock Peer monitoring and credit markets.
\newblock \emph{World Bank Economic Review}, 4(3):351--366, 1990.

\bibitem[UMA(2020)]{uma2020}
UMA Project.
\newblock The Optimistic Oracle.
\newblock UMA Technical Documentation, 2020.
\newblock \url{https://docs.uma.xyz/}

\end{thebibliography}

\end{document}
